\section{Main results}
\label{sec:main}


\subsection{Results for sine-Gordon Y-systems}

\begin{figure}[t]
\centerline{\includegraphics{Nakanishi-Fig1}}
\caption{The diagram $X_n$.}
\label{fig:X}
\end{figure}

With an integer $n \geq 4$, we associate a
diagram $X_n$ in Fig.~\ref{fig:X}.
%For later use, we also attach the properties
%$\circ/\bullet$ and $+/-$ to each vertex of $X_n$.
The diagram $X_n$ should {\em not\/} be regarded as an
ordinary Dynkin diagram,
since the horizontality and verticality of
segments also carry information.
It appeared in \cite{Tateo95}
and encodes the structure of the Y-systems
which we are going to introduce now.
Let $\mathcal{I}_n=\{1,\dots,n+1\}\times \mathbb{Z}$.

\begin{Definition}
\label{def:SGY}
Fix an integer $n \geq 4$.
The sine-Gordon (SG)
Y-system $\mathbb{Y}_{n}(\mathrm{SG})$
is the following system of relations for
a family of variables
 $\{Y_i(u) \mid (i,u)\in \mathcal{I}_n \}$,
\begin{gather}
Y_1\left(u-n+1\right)
Y_1\left(u+n-1\right)
=
\left(\prod_{j=2}^{n-1}(1+Y_j(u-n+j))(1+Y_j(u+n-j))\right)\nonumber\\
\phantom{Y_1\left(u-n+1\right)
Y_1\left(u+n-1\right)
=}{}
 \times (1+Y_n(u))(1+Y_{n+1}(u)),
\nonumber\\
Y_2(u-1) Y_2(u+1)
=
\frac
{1+Y_1(u)}
{1+Y_3(u)^{-1}}
,\nonumber\\
Y_i(u-1)Y_i(u+1)
=
\frac{1}{\prod\limits_{j:j\sim i} (1+Y_j(u)^{-1})},
\qquad i=3,\dots,n+1,\label{eq:Y1}
\end{gather}
where $j\sim i$ means that $j$ is adjacent to $i$ in $X_n$.
\end{Definition}

In \cite{Tateo95}, a more general
family of Y-systems was associated with
a rational parameter $\xi$,
which is related  the coupling constant $\beta$
 of the SG model by \eqref{eq:xi}.
The system \eqref{eq:Y1} corresponds to
the special case
\begin{gather}
\label{eq:case1}
\xi=\frac{n-1}{n}=
\cfrac{1}{ 1 +
\cfrac{1}{n-1}
},
\end{gather}
namely, $F=2$, $n_1=1$, and $n_2=n$ in the notation
in \cite{Tateo95}.
The variable $u$ here is related to the variable $\theta$
in \cite{Tateo95} by $u = (2n/\pi\sqrt{-1})\theta$.
Later in Section~\ref{subsec:back} we explain more about the background
of~\eqref{eq:Y1}.

\begin{Definition}
\label{def:YC}
Let $\EuScript{Y}_{n}(\mathrm{SG})$
be the semif\/ield (Appendix \ref{sec:groc}(i)) with generators
$Y_i(u)$
 $((i,u)\in \mathcal{I}_{n})$
and relations $\mathbb{Y}_{n}(\mathrm{SG})$.
Let $\EuScript{Y}^{\circ}_{n}(\mathrm{SG})$
be the multiplicative subgroup
of $\EuScript{Y}_{n}(\mathrm{SG})$
generated by
$Y_i(u)$, $1+Y_i(u)$
 $((i,u)\in \mathcal{I}_{n})$.
(Here we use the symbol $+$ instead of $\oplus$
for simplicity.)
\end{Definition}
The f\/irst main result of the paper is
the following two theorems conjectured by
\cite{Tateo95}.

\begin{Theorem}[Periodicity]
\label{thm:Yperiod}
The following relations hold in $\EuScript{Y}^{\circ}_{n}(\mathrm{SG})$.
\begin{enumerate}\itemsep=0pt
\item[$(i)$] Half periodicity: $Y_i(u+4n-2)=Y_{\omega(i)}(u)$, where
$\omega$ is an involution of the set $\{1,\dots,n+1\}$
defined by $\omega(n)=n+1$, $\omega(n+1)=n$,
and $\omega(i)=i$ $(i=1,\dots,n-1)$.
\item[$(ii)$] Full periodicity: $Y_i(u+8n-4)=Y_i(u)$.
\end{enumerate}
\end{Theorem}

In our proof of Theorem \ref{thm:Yperiod}
we have  a natural interpretation
of the  half period
\[
4n-2=h(D_n)+2+h(D_{n-1})+2
\]
in terms of the Coxeter number  $h(D_n)=2n-2$ of type $D_n$.

Let $L(x)$ be the {\em Rogers dilogarithm function\/}
 \cite{Lewin81}
\begin{gather*}
%\label{eq:L0}
L(x)=-\frac{1}{2}\int_{0}^x
\left\{ \frac{\log(1-y)}{y}+
\frac{\log y}{1-y}
\right\} dy,
\qquad 0\leq x\leq 1.
\end{gather*}
It satisf\/ies the following relation
\begin{gather}
\label{eq:euler}
L(x)+L(1-x)=\frac{\pi^2}{6},
\qquad 0\leq x\leq 1.
\end{gather}

\begin{Theorem}[Functional dilogarithm identities]
\label{thm:dilog}
Suppose that a family of positive real numbers
$\{ Y_i(u) \mid (i,u)\in \mathcal{I}_n\}$ satisfies
$\mathbb{Y}_n(\mathrm{SG})$.
Then, we have the identities
\begin{gather}
\label{eq:DI}
\frac{6}{\pi^2}
\sum_{
\genfrac{}{}{0pt}{1}
{
(i,u)\in \mathcal{I}_{n}
}
{
0\leq u < 8n-4
}
}
L\left(
\frac{Y_i(u)}{1+Y_i(u)}
\right)
 =8(2n-1),\\
\label{eq:DI'}
\frac{6}{\pi^2}
\sum_{
\genfrac{}{}{0pt}{1}
{
(i,u)\in \mathcal{I}_{n}
}
{
0\leq u < 8n-4
}
}
L\left(
\frac{1}{1+Y_i(u)}
\right)
 =4(n-1)(2n-1).
\end{gather}
\end{Theorem}
Two identities
\eqref{eq:DI} and \eqref{eq:DI'}
are equivalent to each other due to \eqref{eq:euler}.


Using this opportunity, we introduce
another system of relations
 accompanying  $\mathbb{Y}_{n}(\mathrm{SG})$.
\begin{Definition}
Fix an integer $n \geq 4$.
The sine-Gordon (SG)
T-system $\mathbb{T}_{n}(\mathrm{SG})$
is the following system of relations for
a family of variables
 $\{T_i(u) \mid (i,u)\in \mathcal{I}_n \}$,
\begin{gather}
T_1\left(u-n+1\right)
T_1\left(u+n-1\right)
=
T_2(u)+1,\nonumber\\
T_2(u-1) T_2(u+1)
=
T_1(u-n+2)T_1(u+n-2)+T_3(u),
\nonumber\\
T_i(u-1)T_i(u+1)
=
T_1(u-n+i)T_1(u+n-i)
+ \prod_{j:j\sim i} T_j(u),
\qquad i=3,\dots,n-1,
\nonumber\\
T_{n}(u-1)T_{n}(u+1)
=
T_1(u) + T_{n-1}(u),
\nonumber\\
T_{n+1}(u-1)T_{n+1}(u+1)
=
T_1(u) + T_{n-1}(u),\label{eq:T1}
\end{gather}
where $j\sim i$ means that $j$ is adjacent to $i$ in $X_n$.
\end{Definition}

There are two connections between
$\mathbb{Y}_{n}(\mathrm{SG})$ and $\mathbb{T}_{n}(\mathrm{SG})$.
The f\/irst connection is a formal one.
Set
\begin{gather}
\label{eq:di}
d_1=n-1, \qquad d_i=1,\qquad i=2,\dots,n-1,
\end{gather}
and let us  write \eqref{eq:Y1}
in a unif\/ied manner as
\begin{gather}
\label{eq:Yuni}
Y_i\left(u-d_i\right)
Y_i\left(u+d_i\right)
=
\frac{
{\prod\limits_{(j,v)\in \mathcal{I}_{n}}
(1+Y_j(v))^{G_+(j,v;i,u)}
}
}
{\prod\limits_{(j,v)\in \mathcal{I}_{n}}
(1+Y_j(v)^{-1})^{G_-(j,v;i,u)}
}.
\end{gather}
Then, \eqref{eq:T1} is written as
\begin{gather}
\label{eq:Tuni}
T_i\left(u-d_i\right)
T_i\left(u+d_i\right)
=
\prod_{(j,v)\in \mathcal{I}_{n}}
T_j(v)^{G_+(i,u;j,v)}
+
\prod_{(j,v)\in \mathcal{I}_{n}}
T_j(v)^{G_-(i,u;j,v)}.
\end{gather}
Note that we took the `transpositions'
of $G_+$ and $G_-$ in \eqref{eq:Tuni}.
The second connection is an algebraic one.
Suppose that $\{ T_i(u) \mid (i,u)
\in \mathcal{I}_n\}$ satisf\/ies
the T-system $\mathbb{T}_n(\mathrm{SG})$.
Set
\[
Y_i(u)
=
\frac{\prod\limits_{(j,v)\in \mathcal{I}_{n}}
T_j(v)^{G_+(i,u;j,v)}
}
{\prod\limits_{(j,v)\in \mathcal{I}_{n}}
T_j(v)^{G_-(i,u;j,v)}}.
\]
Then, $\{ Y_i(u) \mid (i,u)
\in \mathcal{I}_n\}$ satisf\/ies
the Y-system $\mathbb{Y}_n(\mathrm{SG})$.
One may check the claim directly at this moment
using $\mathbb{T}_n(\mathrm{SG})$
(with some ef\/fort).
Alternatively and better, one can automatically obtain  it from
 \cite[Proposition 3.9]{Fomin07}
once we formulate these systems by a cluster algebra
in the next section.

\begin{Definition}
\label{defn:TC}
Let $\EuScript{T}_{n}(\mathrm{SG})$
be the commutative ring over $\mathbb{Z}$
with identity element,  with gene\-ra\-tors
$T_i(u)^{\pm 1}$ $((i,u)\in \mathcal{I}_{n})$
and relations $\mathbb{T}_{n}(\mathrm{SG})$
together with $T_i(u)T_i(u)^{-1}=1$.
Let $\EuScript{T}^{\circ}_{n}(\mathrm{SG})$
be the subring of $\EuScript{T}_{n}(\mathrm{SG})$
generated by
$T_i(u)$ $((i,u)\in \mathcal{I}_{n})$.
\end{Definition}

The following theorem can be proved simultaneously
with Theorem \ref{thm:Yperiod}.

\begin{Theorem}[Periodicity]
\label{thm:Tperiod}
The following relations hold in $\EuScript{T}^{\circ}_{n}(\mathrm{SG})$.
\begin{enumerate}\itemsep=0pt
\item[$(i)$] Half periodicity: $T_i(u+4n-2)=T_{\omega(i)}(u)$, where
$\omega$ is the  one in Theorem~{\rm \ref{thm:Yperiod}}.
\item[$(ii)$] Full periodicity: $T_i(u+8n-4)=T_i(u)$.
\end{enumerate}
\end{Theorem}

\begin{Remark}
Actually,
 $\mathbb{Y}_n(\mathrm{SG})$ and  $\mathbb{T}_n(\mathrm{SG})$
are also considered for $n=3$,
and they coincide with the Y and T-systems of
type $B_2$ with level $2$ in~\cite{Kuniba94a}.
Theorems~\ref{thm:Yperiod}, \ref{thm:dilog},
and~\ref{thm:Tperiod} remain valid for $n=3$
due to~\cite{Inoue10a}.
\end{Remark}

All the results in this subsection will be
extended to a more general case~\eqref{eq:case2}
later in Section~\ref{sec:extension}.

